\exercisesfor{1}

\begin{enumerate}
\def\labelenumi{\arabic{enumi}.}
\tightlist
\item
  设 G 是有限维连通实或复李群。设 H 和 L 是 G 的李子群，满足 HL = G。证明典范映射 \(G \rightarrow G/H\) 和 \(G \rightarrow G/L\) 定义出解析流形的一个同构
\end{enumerate}

\[
\mathrm{G} / (\mathrm{H} \cap \mathrm{L}) \rightarrow (\mathrm{G} / \mathrm{H}) \times (\mathrm{G} / \mathrm{L}).
\]

\begin{enumerate}
\def\labelenumi{\arabic{enumi}.}
\setcounter{enumi}{1}
\item
  设 \(\mathrm{P} = \{z\in \mathbf{C}|\mathcal{I}(z) > 0\}\)。设 \(\mathbf{G}\) 是由双射 \(z\mapsto az + b\) 在 \(\mathrm{P}(a > 0,b\in \mathbf{R})\) 上组成的群。于是 \(\mathbf{G}\) 在 \(\mathrm{P}\) 上单纯传递地作用，这使我们能够将 \(\mathbf{P}\) 的复解析流形结构传递给 \(\mathbf{G}\)。证明所得结构在 \(\mathbf{G}\) 的左平移下不变，但在右平移下不变。
\item
  设 \(R_{d}\) 是带离散流形结构的实数加法群。令 \(R_{d}\) 按运算律 \((x,y)\mapsto x+y\) 作用在解析流形 R 上。则 \(R_{d}\) 在 R 上传递地作用，但 R 不是 \(R_{d}\) 的李齐性空间。
\item
  设 H 是一个紧的实李群，作用于有限维实流形 V；设 V 属于类 \(C^{r}\)，其中 \(r \in N_{R}\)，且 H 在 V 上的作用属于类 \(C^{r}\)。设 \(p \in V\) 在 H 下不变；H 在线性地作用于 V 在 p 处的切空间 T 上。~证明存在 p 的一个 H-稳定开邻域 U 和一个 \(C^{r}\) -态射 \(f: U \to T\)，使得：
\end{enumerate}

\begin{enumerate}
\def\labelenumi{(\alph{enumi})}
\item
  \(f(p) = 0\)，且 \(f\) 在 \(p\) 处的切映射是恒等映射；
\item
  \(f\) 与 \(\mathbf{H}\) 在 \(\mathbf{U}\) 和 \(\mathbf{T}\) 上的作用相交换。
\end{enumerate}

(先选取 \(f_0\)，使其满足 \((a)\)，再按公式定义 \(f\)

\[
f (x) = \int_ {\mathbf {H}} h. f _ {0} (h ^ {- 1} x) d h,
\]

其中 \(dh\) 是 \(H\) 上总质量为 1 的 Haar 测度。)

推出 H-空间 V 和 T 在 p 的某个邻域内局部同构；换言之，存在 V 在 p 处的坐标系，使得 H 关于这些坐标线性作用（``Bochner 定理``）。

\begin{enumerate}
\def\labelenumi{\arabic{enumi}.}
\setcounter{enumi}{4}
\tightlist
\item
  对超度量域 K 上的流形完成练习 4，假设 H 是一个有限群，其阶 n 满足 \(n.1 \neq 0\) 在 K 中不为 0。
\end{enumerate}

¶ 6. 设 G 是一个实李群，适当地作用于有限维 Hausdorff 实解析流形 V。设 \(p \in V\)，令 H 为 \(p\) 的稳定子群，令 Gp 为 \(p\) 的轨道。H 是紧的，且流形 Gp 可与李齐性空间 G/H 识别。

\begin{enumerate}
\def\labelenumi{(\alph{enumi})}
\item
  证明存在一个穿过 p 的 V 的子流形 S，在 H 下稳定，并且 \(\mathrm{T}_{p}(\mathrm{V})\) 是 \(\mathrm{T}_{p}(\mathrm{G}p)\) 与 \(\mathrm{T}_{p}(\mathrm{S})\) 的直和。（选取 \(\mathrm{T}_{p}(\mathrm{G}p)\) 在 \(\mathrm{T}_{p}(\mathrm{V})\) 中的一个在 H 下稳定的补空间，并应用练习 4。）
\item
  假设 S 如上选取。H 按 \(h.(g,s)=(gh,h^{-1}s)\) 适当地且自由地作用于 G × S；令 \(\mathrm{E}=(\mathrm{G}\times\mathrm{S})/\mathrm{H}\) 为相应的商流形（参见~《微分与解析流形》，R，6.5.1）。映射 \((g,s)\mapsto gs\) 在取商后定义一个态射 \(\rho:E\to V\)。证明 \(\rho\) 与 G 在 E 和 V 上的作用相交换，并且存在一个包含 p、在 H 下稳定的 S 的开子流形 \(S'\)，使得 \(\rho\) 将 \(\mathrm{E}'=(\mathrm{G}\times\mathrm{S}')/\mathrm{H}\) 诱导为轨道 Gp 的一个饱和邻域上的流形同构。
\item
  令 \(N_{p} = T_{p}(V)/T_{p}(Gp)\) 为在 p 处横截于 Gp 的横截空间（《微分与解析流形》，R，5.8.8）；H 作用于 \(N_{p}\)。由 (b) 推出，知道 H 及 H 在 \(N_{p}\) 上的线性表示，就能在 p 的轨道的某个邻域内确定 G-空间 V。
\item
  若 \(x \in \mathbf{N}_p\)，令 \(\mathbf{H}_x\) 为 \(x\) 在 \(\mathbf{H}\) 中的稳定子群。证明存在一个饱和邻域 \(V'\)，它是 \(Gp\) 的邻域，具有如下性质：
\end{enumerate}

对所有 \(p' \in V'\)，存在 \(x \in \mathbb{N}_p\)，使得 \(p'\) 在 \(G\) 中的稳定子群与（在 \(G\) 中）\(H_x\) 共轭。

\begin{enumerate}
\def\labelenumi{(\alph{enumi})}
\setcounter{enumi}{4}
\tightlist
\item
  证明 \(H_{x}\) 在 H 中按共轭仅有有限多个。（对 \(N_{p}=n\) 作归纳，利用 H 在一个由 H 稳定的 n-1 维球面上的作用。）†
\end{enumerate}

\begin{enumerate}
\def\labelenumi{\arabic{enumi}.}
\setcounter{enumi}{6}
\tightlist
\item
  设 E 是 R 上的完备可赋范空间，F 是 E 的闭向量子空间，并且不存在从 E 到 \(\mathbf{F} \times (\mathbf{E}/\mathbf{F})\) 的双连续线性双射。令 X 为解析流形 \(R \times E \times F\)。令 G 为 F 的加法群，它通过映射
\end{enumerate}

\[
(g, (\lambda , e, f)) \mapsto (\lambda , e + g, f + \lambda g)
\]

从 \(\mathbf{G} \times \mathbf{X}\) 到 \(\mathbf{X}\) 的作用映射。对所有 \(x \in \mathbf{X}\)，令 \(\mathbf{R}_x\) 为 \(x\) 的 G-轨道，它是 \(\mathbf{X}\) 的准子流形。于是命题 10 的条件 (a) 得到满足，但不存在具有下列性质的有序对 \((\mathbf{Y}, \pi)\)：\(\mathbf{Y}\) 是解析流形，\(\pi\) 是从 \(\mathbf{X}\) 到 \(\mathbf{Y}\) 的态射，且对所有 \(x \in \mathbf{X}\)，

\[
\mathrm{T} _ {x} (\pi): \mathrm{T} _ {x} (\mathrm{X}) \rightarrow \mathrm{T} _ {\pi (x)} (\mathrm{Y})
\]

是满射的，且核为 \(\mathrm{T}_x(\mathbf{R}_x)\)。

（假设 \((\mathrm{Y},\pi)\) 存在。令 \(\mathbf{H}\) 为 \(\mathbf{Y}\) 在 \(\pi (0,0,0)\) 处的切空间。于是 \(\mathbf{H}\) 同构于 \(\mathbf{R}\times \mathbf{F}\times (\mathrm{E} / \mathrm{F})\)。另一方面，\(\mathbf{H}\) 与 \(\mathrm{T}_{\pi (x)}\mathrm{Y}\) 同构，当 \(x\) 充分接近 \((0,0,0)\) 时，因而同构于 \(\mathbf{R}\times \mathbf{E}\)。）

\begin{enumerate}
\def\labelenumi{\arabic{enumi}.}
\setcounter{enumi}{7}
\item
  令 \(\mathbf{T}\) 带有归一化 Haar 测度，并在 \(\mathcal{L}(\mathrm{L}^2 (\mathbf{T}))\) 上赋予范数拓扑。证明 \(\mathbf{T}\) 在 \(\mathrm{L}^2 (\mathbf{T})\) 上的正则表示不连续。（对所有 \(g\in \mathbf{T}\)（其不同于 \(\epsilon\)），构造 \(f\in \mathrm{L}^2 (\mathbf{T})\)，使得 \(\| f\| = 1\)，\(\| \gamma (g)f - f\| = \sqrt{2}\)。）
\item
  令 I 是由 \((x, \sqrt{2} x) \in \mathbf{R}^2\) 组成的集合，其中 \(-\frac{1}{2} < x < \frac{1}{2}\)。令 U 为 I 在实李群 \(\mathbf{H} = \mathbf{T}^2\) 中的典范像。证明 U 是 H 的李子群芽，但由 U 生成的子群 \(\mathrm{H}'\) 在 H 中稠密且不闭。
\item
  令 \(\mathbf{R}_d\) 为带离散流形结构的群 \(\mathbf{R}\)。令 \(\mathrm{H}\) 为实李群 \(\mathbf{R} \times \mathbf{R}_d\)。令 \(\mathrm{U}\) 为 \(\mathrm{H}\) 中形如 \((x, x)\) 或形如 \((x, -x)\) 的元素集合，其中 \(x \in \mathbf{R}\)。于是 \(\mathrm{U}\) 是 \(\mathrm{H}\) 的李子群芽，\(\mathrm{U}\) 是离散的，\(\mathrm{G}\) 是 \(\mathrm{H}\) 的由 \(\mathrm{U}\) 生成的子群，且等于 \(\mathrm{H}\)；命题 22 的推论在 \(\mathrm{G}\) 上定义的李群结构是离散群结构。
\end{enumerate}

